\section{What the marked theorem now requires}

% BEGIN H70M1
\paragraph{H70M1. Applying the three arithmetic patching inputs on the same chain.}
Fix the data in N69E4, including the original cohomology package
$\mathrm{CohP}$, local reciprocity, one Hida ideal $J$ with its
Taylor--Wiles automorphy system, and a finite chain of admissible
points. N69E1--3 normalize every point, including both blocks at an
exceptional point. Q70C7 chooses the detecting tuple and the
coefficient extensions for the whole chain at once. Apply the
arithmetic places theorem CAP13 to that one tuple. On the resulting
levels N69E4 verifies the relative pro-$\ell$ condition and the
ordered eigenvalue hypotheses. T70B4 and U70O8 then give the two
bounds for each of these very points and levels. Each bound allows
a threshold for each truncation $m$. For fixed $m$ take the maximum
of the finitely many thresholds for the vertices and the two
bounds. The constants fixed before $m$ may also be enlarged over
this finite list. The normalized maps preserve the original kernels.
These are the simultaneous data required by the prime comparison.

R70P1--5 transport the prime comparison at the normalized completions
to the original endpoints and across each edge of the chain. In
particular the dimension argument in R70P4 handles any number of
odd generators. At a genuine two-block vertex the remaining local
work is to construct the source's universal matrix data and its
Selmer presentation on that completion. This includes passing the
bounds and auxiliary eigenvalues just obtained to the limiting
group in the form needed by the presentation. For a full-block
vertex R70P6 gives the single-block comparison directly. Given
the indicated two-block comparison at every required vertex, the
preceding argument proves $\mathrm{Raw}$ on the chains of N69E4
and hence both prime tests in N69E5. This proves the original
$\mathrm{PfromAut}$ under the indicated local construction and
the stated arithmetic foundations.
% END H70M1

% BEGIN H70M2
\paragraph{H70M2. Finiteness and descent use one mark on each host.}
Assume the cohomological inputs for the original marked hosts and
patching data, and the original Hida inputs $H'$ of HIA2.
Suppose the local construction of H70M1 is available on all
the required two-block completions. HIA2 gives one Hida ideal
with its automorphy system on each arithmetic setup. KX3
gives the local reciprocity coordinates. H70M1 and N69E5
provide the two prime tests on that same ideal. HIA3 gives the
unitary modularity-ideal assertion. By L69D3 the chosen link
therefore has weight-zero finiteness and upper-field automorphy
in the required two ranks. D70W5 now supplies the polarized
solvable descent needed in L69D4. In particular the inverse
cyclotomic twist in rank three is attached over the Gaussian
field with its one real window center for each attached cusp.

Apply M65A4 on the original host. The mark is chosen before the
universal ring or the ordinary point. The finite conclusion and
the automorphic conclusion hold at this same mark, as required
by the quantifiers of $\mathrm{ThmS6}_s$. This conditional proof
locates the remaining work in the two-block comparison and the
original Hida and cohomological inputs. Its conclusion has no
additional solvable descent premise.
% END H70M2
